\documentclass[11pt]{article}
\usepackage[margin=1in]{geometry}
\usepackage{amsmath,amssymb}
\usepackage{enumitem}
\usepackage{hyperref}
\title{Response to Reviewers (Paper ID 296)\\
\large Integrating Adaptive Neuro-Fuzzy Inference Systems and Large Language Models for Arabic Dialect Processing}
\author{}
\date{}

\begin{document}
\maketitle
\vspace{-1.5em}

\noindent\emph{Note: Equation and table numbers below (e.g., Eq.~1, Table~I) refer to the revised manuscript.}

\section*{Response to Reviewers (Paper ID 296)}
We thank the reviewers and editor for the careful and constructive assessment of our manuscript, ``Integrating Adaptive Neuro-Fuzzy Inference Systems and Large Language Models for Arabic Dialect Processing'' (Paper ID 296). Below we reproduce each reviewer's comment in full and give a point-by-point response describing the corresponding changes made in this revised manuscript.

\section*{Reviewer 1 (Confidence: High)}
\textit{``The paper addresses an interesting problem in Arabic NLP. Treating dialect boundaries as potentially gradual rather than strictly categorical is a reasonable motivation, and the use of parallel dialogue scenarios provides a controlled setting for examining representation differences across dialects. I initially expected the paper, based on its title and abstract, to experimentally evaluate an integrated ANFIS--LLM system. However, this is not what the current study actually does. The implemented experimental component consists of Sentence Transformer embeddings followed by PCA, centroid-distance analysis, and t-SNE visualization. The ANFIS component is proposed architecturally but is neither trained nor evaluated. The authors themselves acknowledge this explicitly in the limitations section: `no end-to-end implementation or experimental evaluation of the proposed hybrid model is presented.'''}

\textbf{Response:} We thank the reviewer for this precise and entirely fair criticism, which identified a genuine gap between our original manuscript's framing and its actual empirical content at the time of submission. Rather than revise the title or abstract to lower expectations, we chose to close this gap directly by implementing, training, and evaluating a real ANFIS layer, and we have substantially revised the manuscript accordingly:

\begin{enumerate}
\item \textbf{Methodology (Section~II-A, ``Training and Evaluation of the ANFIS Layer,'' new).} We now describe a concrete zero-order Sugeno-type ANFIS classifier. Centroids and per-class, per-dimension (anisotropic) variances (Eq.~3) are estimated exclusively from an 80\% training split (stratified by dialect, random seed 42); the remaining 20\% is held out strictly for evaluation. The PCA projection itself is also fit only on the training partition and then applied to the held-out test partition, so that no test information leaks into any stage of model fitting. We further add a learned defuzzification (consequent) layer (Eq.~5), completing the canonical five-layer Sugeno ANFIS architecture rather than deciding directly from normalized membership degrees, and we describe a $k$-nearest-neighbor ($k$=5) baseline fit on the identical split for comparison.

\item \textbf{Results (Section~III-C, ``ANFIS Dialect-Membership Evaluation,'' new).} We report held-out test performance: an overall accuracy of 85.33\% (600 held-out test samples), compared with a 33.3\% chance baseline and a 79.00\% $k$-NN baseline computed on the same embeddings (Table~III), with the ANFIS layer now exceeding the $k$-NN baseline by a clear margin. We include a full per-dialect precision/recall/F1 breakdown and confusion matrix (Table~IV), a controlled ablation isolating the effect of the anisotropic premise versus the defuzzification layer (Table~II), and concrete qualitative examples showing the fuzzy membership vector $\mu_d$ behaving as intended: near-unanimous membership for an unambiguous utterance, versus genuinely graded, distributed membership for a more ambiguous one, with the correct dialect recovered in both cases.

\item \textbf{Limitations (Section~IV-A, revised).} We removed the sentence quoted by the reviewer (``no end-to-end implementation or experimental evaluation of the proposed hybrid model is presented''), since it is no longer accurate. This subsection is retitled ``Residual Classification Errors and Scope of the Embedding-Model Swap'' and now honestly discusses the roughly 15\% of held-out examples still misclassified, and reports that the defuzzification layer's own contribution to accuracy is within noise for $n=600$---we present it as completing the architecture, not as an accuracy-improving step, since we do not want to overstate a change that the ablation shows is not reliably positive.

\item \textbf{Abstract and Conclusion (revised).} Both now state the trained/evaluated result explicitly (85.33\% accuracy, above chance level and above the $k$-NN baseline, with balanced per-dialect recall) so that the paper's framing matches its actual empirical content, directly resolving the title/abstract-versus-content mismatch the reviewer identified.

\item \textbf{Beyond the minimum fix.} Having closed the specific gap the reviewer identified, we went further: we replaced the general-purpose sentence embedding used for this classification sub-task with a dialect-specialized encoder (MARBERT), which is what drove accuracy from the mid-40s to the mid-80s (an ablation in Table~II isolates this and the other architectural changes individually), and we instantiated a working closed-loop ANFIS--LLM pipeline in which the trained ANFIS layer's membership degrees steer an LLM decoder toward dialect-appropriate generation, illustrated qualitatively in the revised Conclusion and documented in the supplementary material.
\end{enumerate}

We believe these changes directly and substantively address the reviewer's central concern: the ANFIS component is no longer merely proposed architecturally, but is trained on a genuine train/test split and evaluated on held-out data with reported accuracy, a confusion matrix, an ablation, and a baseline comparison that the trained model now clearly exceeds. We are grateful to the reviewer for pushing us to complete this work rather than simply reframe our claims, and we believe the paper is substantially stronger as a result.

\section*{Reviewer 2 (Confidence: High)}
\textit{``I found the dataset design interesting... My main concern is with the interpretation of the embedding geometry. Table I reports the following utterance centroid distances... These values show differences between centroid locations, but centroid distance alone does not establish whether the dialect distributions are statistically distinguishable. The t-SNE plots on page 3 are visually informative, but I would be careful with statements such as `clearly distinguishable' or `minimal overlap.' t-SNE is primarily a visualization technique and does not preserve global geometry reliably enough to serve as quantitative evidence of separability... More importantly, there is a mismatch between the paper's fuzzy-membership motivation and its empirical evidence... the corpus contains three manually created parallel versions with hard gold labels. There are no actual mixed-dialect or code-switched samples for which graded membership has been annotated or independently validated... I would suggest a relatively simple experiment... train a baseline dialect classifier directly on the embeddings using cross-validation and report macro-F1 and confusion matrices... If the authors implement ANFIS, its membership outputs could then be compared with classifier uncertainty or, preferably, human judgments on genuinely mixed-dialect examples.''}

\textbf{Response:} We thank the reviewer for a precise and constructive critique that pushed us to separate what our data can and cannot support. We address each point directly:

\begin{enumerate}
\item \textbf{Centroid distance is not statistical evidence of separability.} We agree, and no longer rely on it as such. Section~III now explicitly states that Sections~III-A and III-B (centroid distances, t-SNE) are a descriptive, unsupervised analysis over the full corpus, appropriate for exploratory characterization but not predictive validation, while Section~III-C's held-out, supervised ANFIS evaluation is what we offer as actual quantitative evidence of separability. In addition, we now back the centroid-distance claim itself with a proper statistical test rather than raw distance values alone: Section~III-A reports a permutation test in which the dialect labels are randomly shuffled 1,000 times and the average pairwise centroid distance is recomputed under each shuffle. The observed distance under the true labels (0.389) far exceeds the distribution obtained under random shuffling (mean 0.117, std 0.011, max 0.157), $p=0.001$, giving genuine statistical evidence that the centroid separation is not due to chance.

\item \textbf{Overclaiming language around t-SNE (``clearly distinguishable,'' ``minimal overlap'').} We have removed this language. Section~III-B now states plainly that t-SNE does not reliably preserve global distances, that some mixing among all three dialects persists even where GUL appears comparatively separated, and that the plots are illustrative rather than quantitative evidence, redirecting the reader to the held-out classification results in Section~III-C for the actual evidence of separability. We also add the silhouette score for the three dialect labels in this same 50-dimensional PCA space, which is $-0.0189$---essentially zero, confirming quantitatively (not just visually) that individual utterances are heavily overlapping at the point level even though their class centroids are significantly separated (per the permutation test above); we present these two statistics together as the honest picture of this unsupervised space, which is precisely what motivates treating dialect boundaries as fuzzy rather than hard categories.

\item \textbf{The suggested experiment: a baseline classifier with cross-validation, macro-F1, and confusion matrices.} We implemented this suggestion in full, including the cross-validation the reviewer specifically requested. Section~III-C reports a zero-order Sugeno ANFIS classifier trained and evaluated on a stratified 80/20 held-out split, with a full confusion matrix (Table~IV), per-dialect precision/recall/F1 and a macro-averaged F1-score of 0.8533 (Table~III), and a $k$-nearest-neighbor baseline on the identical embeddings and split (0.7900), which the trained ANFIS layer now clearly exceeds. We additionally ran 5-fold stratified cross-validation of the identical model (PCA refit within each fold to prevent leakage), obtaining a mean accuracy of 0.8499 (std 0.0132) and a mean macro-F1 of 0.8498 (std 0.0134) across folds, closely matching the single-split result and confirming it is not an artifact of one favorable partition (Section~III-C).

\item \textbf{Mismatch between the fuzzy-membership motivation and the empirical evidence; no genuinely mixed-dialect, independently annotated data.} We agree this gap cannot be fully closed without genuinely annotated mixed-dialect data, but we no longer leave it entirely untested. We built a synthetic mixed-dialect probe: for held-out test-split scenarios (present in all three dialects, since the corpus is parallel), we spliced the first half of one dialect's utterance with the second half of another dialect's utterance of the same scenario, producing 264 synthetic mixed sentences with a known true pair of source dialects. The trained ANFIS layer's average top-membership confidence dropped from 0.868 on pure sentences to 0.805 on these synthetic mixes, and its top-two membership dialects matched the true spliced-together pair in 58.0\% of cases---well above the 33\% chance level for guessing among the three possible dialect pairs. This is a real, run, quantitative signal that the membership degrees track genuine mixing rather than arbitrary noise, and we report it as such. We are careful, however, not to overstate it: word-splicing is a crude synthetic approximation of mixing, not natural code-switching by a bilingual speaker, no human verified the resulting sentences' naturalness, and 58\% top-2 agreement, though well above chance, is far from perfect. We therefore retain this as an explicit limitation (Section~IV-A, ``Fourth''), now with partial supporting evidence rather than none, and continue to recommend a genuinely annotated mixed-dialect dataset, compared against classifier confidence or human judgment exactly as the reviewer suggests, as the definitive future-work direction.
\end{enumerate}

We are grateful for a review that distinguished sharply between what our analysis illustrates and what it actually proves; we believe the manuscript is now considerably more honest about that distinction, and correspondingly stronger.

\section*{Reviewer 3 (Confidence: High; Recommendation: Reject)}
\textit{``...Figure 1 on page 3 presents a complete pipeline consisting of user input, LLM/encoder, PCA embedding space, ANFIS layer, LLM decoder/response generator, and dialect-adapted response... None of these ANFIS or response-generation stages is experimentally implemented in the reported study... the implemented experiment is essentially: Sentence Transformer $\to$ PCA $\to$ centroid distances $\to$ t-SNE. All of these are standard methods, and the paper does not introduce a new embedding model, dimensionality-reduction algorithm, dialect classifier, fuzzy inference mechanism, or response-generation experiment. I also question the statement that the findings show that `multilingual Transformers can effectively encode the structure of Arabic dialects.' The dataset consists of parallel versions of the same 999 scenarios... there is no external dataset or naturally occurring corpus to demonstrate generalization. There is also no train/test protocol because the primary analysis is unsupervised... this should not be presented as validation of a proposed dialogue architecture... In my view, however, this admission also identifies why the manuscript is premature as an experimental paper.''}

\textbf{Response:} We thank the reviewer for a direct and demanding assessment, and we take each objection in turn, since several of them describe exactly the manuscript the reviewer saw at the time of the original review, which we have since substantively changed rather than merely reframed.

\begin{enumerate}
\item \textbf{``None of these ANFIS or response-generation stages is experimentally implemented.''} This was accurate of the original submission and is no longer accurate. The ANFIS stage is now a trained, held-out-evaluated zero-order Sugeno classifier (Section~II-A, anisotropic Gaussian premise plus a learned defuzzification/consequent layer, Eqs.~3--5), reporting 85.33\% accuracy, a full confusion matrix, macro-F1, and a $k$-NN baseline (Section~III-C). The response-generation stage from Fig.~1 is now also real: the trained ANFIS layer's membership degrees are converted into a dialect-aware steering prompt that conditions an LLM decoder, demonstrated on held-out utterances and in ENOSH, our live interactive demo (Fig.~3). We are careful not to overstate this second piece: as noted in Section~IV-A (``Third''), the generation demo is illustrative rather than systematically evaluated with automatic or human quality metrics across the full test set, and that systematic evaluation is left to future work. But the reviewer's specific claim---that these stages are not experimentally implemented at all---no longer holds.

\item \textbf{``The implemented experiment is essentially Sentence Transformer $\to$ PCA $\to$ centroid distances $\to$ t-SNE... no new... dialect classifier, fuzzy inference mechanism, or response-generation experiment.''} This description matched the original manuscript's experimental content. It no longer does: the manuscript now additionally contributes (i) a trained anisotropic-premise, defuzzification-completed ANFIS classifier, validated against an isotropic-premise ablation and a $k$-NN baseline on identical embeddings (Table~II), which is a genuine fuzzy-inference mechanism, not merely a proposed one; (ii) a dialect-specialized embedding swap (MARBERT) for the classification sub-task, whose contribution is isolated by that same ablation rather than asserted; and (iii) a working closed-loop generation experiment (Fig.~3). We agree that the exploratory centroid/t-SNE analysis (Sections~III-A--III-B) by itself would not constitute a sufficient experimental contribution, and no longer ask it to carry that weight; Section~III now says so explicitly.

\item \textbf{Generalization claim (``multilingual Transformers can effectively encode the structure of Arabic dialects'') given a parallel, topic-matched dataset with no external corpus.} We agree this claim was overstated, and have narrowed it. The Abstract and Conclusion now explicitly scope this finding to the parallel, topic-matched corpus used here, state plainly that we do not claim generalization to naturally occurring, topic-unmatched Arabic text, and name evaluation on an external corpus as future work.

\item \textbf{No train/test protocol; the same full dataset is used to construct PCA representations and to draw conclusions, which should not be presented as validation.} We agree with this characterization of Sections~III-A--III-B and no longer present them as validation. Section~III now opens by explicitly separating the two evaluation regimes: the unsupervised, whole-corpus centroid/t-SNE analysis is descriptive and exploratory only, while the ANFIS classifier in Section~III-C is fit exclusively on an 80\% training partition and evaluated strictly on the remaining held-out 20\%---a genuine train/test protocol, distinct from, and not conflated with, the exploratory analysis.

\item \textbf{Overall assessment: the manuscript was premature as an experimental paper.} We accept that this was a fair assessment of the version reviewed. Our response has been to close that gap with real experiments rather than to argue the original framing was adequate: a trained and held-out-evaluated ANFIS classifier with an ablation and baseline, a real closed-loop ANFIS--LLM generation pipeline, and explicit, honest scoping of every claim we can and cannot currently support (Section~IV). We respectfully ask the reviewer to reconsider the manuscript in light of these substantive additions, which we believe address the specific basis for the original recommendation rather than merely reframing around it.
\end{enumerate}

We thank all three reviewers for a rigorous first round; the manuscript is substantially more honest and considerably more empirically grounded as a result.

\end{document}
